\documentclass[10pt,a4paper]{amsart}
\usepackage[margin=19mm]{geometry}
\usepackage{amsmath,amssymb,mathtools}
\usepackage[scr=rsfs,cal=euler]{mathalfa}
\usepackage{xfrac}
\usepackage[unicode,hidelinks,bookmarks=false]{hyperref}
\newcommand{\ie}{i.e.\ }
\newcommand{\cf}{cf.\ }
\newcommand{\resp}{respectively\ }
\newcommand{\Subsection}{\subsection}
\providecommand{\nobreakdash}{\nobreak\hbox{-}}
\setlength{\emergencystretch}{2em}
\multlinegap0pt
\usepackage{bm}
\usepackage{bbm}
\usepackage{dsfont}
\usepackage{xspace}
\usepackage{cancel}
\usepackage{mathtools}
\usepackage{cleveref}
\usepackage{amsthm}
\makeatletter
\@ifundefined{assumption}{\newtheorem{assumption}{Assumption}}{}
\@ifundefined{lemma}{\newtheorem{lemma}{Lemma}}{}
\makeatother
\newcommand{\Ctau}{C_\tau}
\newcommand{\R}{\mathbb{R}}
\newcommand{\abs}[1]{\left\lvert#1\right\rvert}
\newcommand{\normbig}[1]{\big\lVert#1\big\rVert}
\newcommand{\norm}[1]{\left\lVert#1\right\rVert}
\newcommand{\sspa}{\mathbb{S}}
\newcommand{\tmix}{\tau}
\title[Projection-based Lyapunov method for fully heterogeneous wea]{Projection-based Lyapunov method for fully heterogeneous weakly-coupled MDPs}
\author{Xiangcheng Zhang, Yige Hong, Weina Wang}
\date{Source: arXiv:2502.06072v6 (CC BY 4.0)}
\begin{document}
\maketitle

\noindent\emph{Source and licence.} Projection-based Lyapunov method for fully heterogeneous weakly-coupled MDPs. Version arXiv:2502.06072v6 (CC BY 4.0), distributed under the Creative Commons Attribution 4.0 International licence, as recorded in the arXiv licence metadata for this exact version and shown on the abstract page https://arxiv.org/abs/2502.06072v6. Full rights notice: licenses/arxiv-2502.06072-CC-BY-4.0.txt.\par\smallskip

\noindent\textit{Editorial scope.} Counted unit: the Lemma whose source label is \texttt{lem:cgamma} in the article (source0.tex lines 1804--1814, source label \texttt{lem:cgamma}), delivered together with the article's own complete original proof of it (source0.tex lines 1819--1869, one \texttt{proof} environment of 51 lines). No mathematics is written, repaired or re-derived: statement and proof are the article's own text line for line. The proof is detached from the statement in the source: the article states the result at lines 1804--1814 and prints its proof at lines 1819--1869, and the proof environment's own header names the statement, so the association is the article's own. Required context retained uncounted: ass:unichain (lines 1179--1185). Every definition, display and label of those lines is retained unchanged. Omitted from this package and not redistributed: 1--1178, 1186--1803, 1815--1818, 1870--2656. Nothing inside a retained excerpt is omitted or altered: every retained body is delivered exactly as the article's lines are written, so no omission receipt of any kind accompanies this package. Citations: the retained excerpts carry no citation command at all, so no bibliography entry of the article is transcribed and no reference is redistributed. Reference adaptation: none was needed and none was made; every reference inside the delivered unit resolves inside it, so no unresolved reference or citation is produced. Printed slips kept literal: no printed word, symbol, hypothesis or number of the counted statement or of its proof was altered. Layout and class: the article's own class is \texttt{article} with packages including geometry, natbib, amsmath, amsfonts, amssymb, amsthm, inputenc, fontenc, hyperref, url, booktabs, microtype, xcolor, times; the wrapper is the lane's pinned \texttt{amsart} editorial wrapper at 10pt on A4, which restates the theorem-like environments the retained text needs and adds the article's own macro definitions that the retained text needs (\texttt{\textbackslash Ctau}, \texttt{\textbackslash R}, \texttt{\textbackslash abs}, \texttt{\textbackslash norm}, \texttt{\textbackslash normbig}, \texttt{\textbackslash sspa}, \texttt{\textbackslash tmix}), with extra wrapper packages (bm, bbm, dsfont, xspace, cancel, mathtools, cleveref). The delivered numbering is the unit's own. Assets: no retained span contains a figure environment, a graphics-inclusion command, a PDF-page inclusion, a table or a TikZ picture; no image file is copied into this package, the package declares no asset dependency and no image file is redistributed with it. Licence: version arXiv:2502.06072v6 of this article is distributed under the Creative Commons Attribution 4.0 International licence, as recorded in the arXiv licence metadata for this exact version and shown on the abstract page https://arxiv.org/abs/2502.06072v6. A full rights notice is delivered at licenses/arxiv-2502.06072-CC-BY-4.0.txt. Characters: every retained excerpt is pure ASCII, so the delivered engine, XeLaTeX with its default Latin Modern text font, reports no missing character.
\begin{assumption}\label{ass:unichain}
    For each arm $i\in\mathbb{N}_+$, the induced Markov chain under the optimal single-armed policy $\bar{\pi}_i^*$ is an aperiodic unichain. 
    Furthermore, the mixing times of these Markov chains have a uniform upper bound; i.e., there exists a positive $\tau$ such that for all $i\in\mathbb{N}_+$, 
    \begin{equation}
        \tmix_i \leq \tau.
    \end{equation}
\end{assumption}

\begin{lemma}\label{lem:cgamma}
    Suppose \Cref{ass:unichain} holds, and let
    \[
        \gamma = \exp\Big(-\frac{1}{2\tau}\Big) \qquad \Ctau = \frac{4e}{1-1/\sqrt{e}} \tau.
    \]
    Then we have
    \begin{align*}
        \sum_{\ell=0}^{\infty} \max_{i\in[N]}\frac{\norm{(P_i-\Xi_i)^\ell}_\infty}{\gamma^\ell} \leq \Ctau. 
\end{align*}
 Here, recall that $\norm{A}_\infty \triangleq \max_{s\in\sspa} \sum_{s'\in\sspa}\abs{A(s,s')}$ for any matrix $A\in \R^{\sspa\times\sspa}$. 
\end{lemma}

\begin{proof}[Proof of \Cref{lem:cgamma}]
    For each $\ell \in \mathbb{N}$, let $\ell_0 = \ell - \tau \lfloor \ell/\tau\rfloor$. 
    By properties of the operator norm $\norm{\cdot}_\infty$, 
    \begin{equation}
        \label{eq:pf-cgamma-bound:int-divide}
        \norm{(P_i - \Xi_i)^\ell}_\infty \leq \norm{(P_i - \Xi_i)^{\ell_0}}_\infty \norm{(P_i - \Xi_i)^\tau}_\infty^{\lfloor \ell/\tau \rfloor}. 
    \end{equation}
    Because each row of $\Xi_i$ is $\mu_i^*$, it is easy to verify that $(P_i - \Xi_i)^{\ell'} = P_i^{\ell'} - \Xi_i$ for any $\ell' \in \mathbb{N}$. Consequently, 
    \begin{align*}
        \norm{(P_i - \Xi_i)^{\ell_0}}_\infty
        &= \normbig{P_i^{\ell_0} - \Xi_i}_\infty \\
        &= \max_{s\in\sspa} \sum_{s'\in\sspa} \abs{P_i^{\ell_0}(s,s') - \mu_i^*(s')}  \\
        &\leq \max_{s\in\sspa} \sum_{s'\in\sspa} \big(P_i^{\ell_0}(s,s') + \mu_i^*(s')\big) \\
        &= 2.
    \end{align*}
    Moreover, by \Cref{ass:unichain}, it is not hard to see that 
    \[
        \max_{s\in\sspa} \norm{P_i^\tau(s,\cdot) - \mu_i^*(\cdot)}_1 \leq \frac{1}{e},
    \]
    so we have
    \begin{align}
        \nonumber
        \norm{(P_i - \Xi_i)^\tau}_\infty 
        &= \max_{s\in\sspa} \sum_{s'\in\sspa} \abs{P_i^\tau(s,s') - \mu_i^*(s')} \\
        \nonumber
        &= \max_{s\in\sspa} \norm{P_i^\tau(s,\cdot) - \mu_i^*(\cdot)}_1 \\
        \nonumber
        &\leq \frac{1}{e}.
    \end{align}
    Substituting the bounds on $\norm{(P_i - \Xi_i)^\tau}_\infty$ and $\norm{(P_i - \Xi_i)^{\ell_0}}_\infty$ back to \eqref{eq:pf-cgamma-bound:int-divide}, we get 
    \begin{align}
        \nonumber
        \norm{(P_i - \Xi_i)^\ell}_\infty
        &\leq 2\exp\Big(-\Big\lfloor \frac{\ell}{\tau} \Big\rfloor\Big) \\
        &\leq 2\exp\Big(-\frac{\ell}{\tau} + 1\Big).
    \end{align}

    Therefore, 
    \begin{align}
        \nonumber
        \sum_{\ell=0}^{\infty} \max_{i\in[N]}\frac{\norm{(P_i-\Xi_i)^\ell}_\infty}{\gamma^\ell}
        &\leq \sum_{\ell=0}^{\infty} 2\exp\Big(-\frac{\ell}{\tau} + 1\Big) \gamma^{-\ell} \\
        \nonumber
        &= \sum_{\ell=0}^{\infty} 2\exp\Big(-\frac{\ell}{2\tau} + 1\Big) \\
        &= \frac{2e}{1 - \exp(-1/(2\tau))}. 
    \end{align}
    Observe that $\tau \geq 1$, $0 < 1/(2\tau) \leq 1/2$, and $\exp(-a) \leq 1 - (1-1/\sqrt{e}) a$ for any $0 < a \leq 1/2$, so
    \[
        \frac{2e}{1 - \exp(-1/(2\tau))} \leq \frac{4e}{1-1/\sqrt{e}} \tau = \Ctau. \qedhere
    \]
\end{proof}

\end{document}
